\documentclass[11pt,letterpaper]{article}
\usepackage[margin=0.85in,top=0.75in,bottom=0.75in,headheight=14pt]{geometry}
\usepackage[T1]{fontenc}
\usepackage{newtxtext,newtxmath,microtype}
\usepackage{booktabs,tabularx,amsmath,xcolor}
\usepackage{pgfplots}
\pgfplotsset{compat=1.18}
\usepackage{listings,fancyhdr,titlesec}
\usepackage[hidelinks]{hyperref}
\definecolor{ink}{HTML}{202A35}
\definecolor{muted}{HTML}{616973}
\definecolor{rulegray}{HTML}{C5CBD1}
\definecolor{navy}{HTML}{334E68}
\setlength{\parindent}{0pt}
\setlength{\parskip}{6pt}
\linespread{1.04}
\titleformat{\section}{\normalsize\bfseries\color{ink}}{\thesection}{0.65em}{}
\titlespacing*{\section}{0pt}{13pt}{5pt}
\pagestyle{fancy}\fancyhf{}
\fancyhead[L]{\footnotesize\color{muted}RESEARCH METHODS IN FINANCE}
\fancyhead[R]{\footnotesize\color{muted}WEEK 02}
\fancyfoot[L]{\footnotesize\color{muted}Olin Business School\enspace /\enspace Fall 2026}
\fancyfoot[R]{\footnotesize\thepage}
\renewcommand{\headrulewidth}{0pt}
\renewcommand{\footrulewidth}{0pt}
\lstset{basicstyle=\ttfamily\fontsize{8}{10}\selectfont,breaklines=true,columns=fullflexible,keepspaces=true,showstringspaces=false,numbers=left,numberstyle=\tiny\color{muted},numbersep=10pt,xleftmargin=18pt,frame=none,aboveskip=8pt,belowskip=0pt}
\hypersetup{pdftitle={Operating Profitability and Portfolio Excess Returns},pdfauthor={},pdfsubject={Week 2 | Stata Group Assignment 1}}
\begin{document}
{\small\color{muted}STATA GROUP ASSIGNMENT 1}\par
\vspace{3pt}
{\fontsize{24}{28}\selectfont\bfseries\color{ink}Operating Profitability\\and Portfolio Excess Returns}\par
\vspace{4pt}
{\small Value-weighted deciles\enspace\textbar\enspace July 1963--July 2026\enspace\textbar\enspace 757 months}
\vspace{7pt}\hrule height 0.5pt\vspace{5pt}

\section{Data and portfolio construction}
We use Kenneth French's ten value-weighted portfolios formed on operating profitability (OP), ordered from the lowest-profitability decile (1) to the highest (10). Portfolios are formed each June using NYSE breakpoints. OP is prior-year revenue less cost of goods sold, interest expense, and selling, general and administrative expenses, divided by book equity.

Monthly returns are matched by calendar month to the risk-free return (RF) in the Fama/French three-factor file. The common sample contains 757 complete observations per decile. The files were downloaded on September 15, 2026 and use the July 2026 CRSP vintage; only the monthly value-weighted return block is selected.

\section{Excess returns and inference}
For portfolio $d$, monthly excess return is $x_{d,t}=R_{d,t}-R_{f,t}$, measured in percentage points. We test $H_0:\mu_d=0$ against $H_1:\mu_d\ne0$ using
\begin{equation*}
 t_d=\frac{\bar{x}_d}{s_d/\sqrt{N}},\qquad
 s_d^2=\frac{1}{N-1}\sum_{t=1}^{N}(x_{d,t}-\bar{x}_d)^2,\qquad
 p_d=2\Pr\!\left(T_{756}\geq |t_d|\right).
\end{equation*}
\vspace{-6pt}
\section{Results}
{\small\bfseries Table 1.\enspace Monthly excess returns by operating-profitability decile}\par
\vspace{2pt}
\begingroup
\small\renewcommand{\arraystretch}{1.27}
\begin{tabular*}{\textwidth}{@{\extracolsep{\fill}}lrrrr@{}}
\toprule
OP decile & Mean (\%) & Std. dev. (\%) & $t$-statistic & $p$-value\\
\midrule
\csname @@input\endcsname results.tex 
\bottomrule
\end{tabular*}
\endgroup
\vspace{4pt}
{\footnotesize\color{muted}Notes: $N=757$ for every decile. Means are arithmetic monthly averages; standard deviations use $N-1$. Tests are two-sided with 756 degrees of freedom. P-values are unadjusted. Source: Kenneth French Data Library.}\par
\vspace{4pt}
\textbf{Main finding.} Nine of ten deciles have positive mean excess returns significant at 5\%. The lowest-profitability decile does not reject the zero-mean null ($p=0.0672$).

\newpage
\section{Interpretation}
Mean excess returns range from 0.4394\% to 0.7272\% per month. Deciles 2--10 reject the zero-mean null at both 5\% and 1\%. Decile 1 is significant at 10\%, but not at 5\% ($t=1.8327$).

Higher-profitability portfolios generally have larger mean excess returns, although the pattern is not monotonic. Decile 8 has the highest mean (0.7272\%); decile 10 averages 0.6910\%. The lowest-profitability portfolio is also the most volatile: its monthly standard deviation is 6.5959\%, compared with 4.6210\% for decile 10.

The difference between the decile-10 and decile-1 sample means is approximately 0.252 percentage points per month. Individual portfolio tests do not establish the statistical significance of this difference. Standard errors do not adjust for serial correlation, and p-values are not corrected for multiple testing. These results do not establish risk-adjusted abnormal performance or a causal effect of profitability.

\section{Monthly return variation}
{\small\bfseries Figure 1.\enspace Highest-profitability decile: monthly excess return}\par
{\footnotesize\color{muted}Value-weighted OP decile 10\enspace /\enspace July 1963--July 2026}
\begin{center}
\begin{tikzpicture}
\begin{axis}[width=\textwidth,height=2.8in,scale only axis=false,
 xmin=1963.5,xmax=2026.583,ymin=-30,ymax=20,
 xtick={1970,1980,1990,2000,2010,2020},xticklabel style={/pgf/number format/1000 sep=},
 ytick={-30,-20,-10,0,10,20},
 ylabel={Excess return (\% per month)},
 tick label style={font=\footnotesize,color=muted},label style={font=\footnotesize,color=muted},
 axis lines=left,axis line style={rulegray},tick style={rulegray},
 ymajorgrids=true,grid style={rulegray!45,line width=0.25pt},
 clip=true]
\addplot[muted,line width=0.4pt] coordinates {(1963.5,0) (2026.583,0)};
\addplot[navy,line width=0.45pt,no marks] table[x=year,y=excess]{decile10.dat};
\end{axis}
\end{tikzpicture}
\end{center}
\vspace{-8pt}
{\footnotesize\color{muted}Source: Kenneth French Data Library; portfolio return minus RF. The figure uses the same monthly observations as the Stata analysis.}\par
Despite its positive sample mean, decile 10 experiences substantial month-to-month fluctuations. Its largest negative excess return occurs in October 1987 ($-25.17\%$). The figure shows monthly returns, not cumulative investment performance.

\section*{Sources and reproducibility}
{\footnotesize
\href{https://mba.tuck.dartmouth.edu/pages/faculty/ken.french/data_library.html}{Kenneth French Data Library}; \href{https://mba.tuck.dartmouth.edu/pages/faculty/ken.french/Data_Library/det_port_form_op.html}{operating-profitability portfolio definitions}. RF uses Ibbotson one-month Treasury bill data through May 2024 and the ICE BofA US 1-Month Treasury Bill Index thereafter. The library's current CIZ-based research return series is used.

The \texttt{week-02} folder includes frozen source files, checksums, code and outputs. From the repository root in Stata, run \texttt{do run-week-02.do}. The appendix reproduces the analysis do-file. The report's editable LaTeX source is in \texttt{week-02/report/}.}

\newpage
{\small\color{muted}APPENDIX A}\par
{\Large\bfseries Stata code}\par
{\small\color{muted}Data preparation and monthly excess returns}
\lstinputlisting[firstline=1,lastline=60]{../assignment1.do}
\newpage
{\small\color{muted}APPENDIX A\enspace /\enspace CONTINUED}\par
{\Large\bfseries Stata code}\par
{\small\color{muted}Hypothesis tests and Stata figure export}
\lstinputlisting[firstline=61,firstnumber=61]{../assignment1.do}
\end{document}
